% ============================================================
% Thesis: Neighbor-Interpolated Point Cloud Densification
%         for 3D Gaussian Splatting Initialization
% Author: ABDUGANIEV ABDULVOSIT
% Seoul National University of Science and Technology
% MSc Integrated IT Engineering, 2026
% ============================================================

\documentclass[12pt, a4paper]{article}

% --- Packages ---
\usepackage[T1]{fontenc}
\usepackage[utf8]{inputenc}

\usepackage[margin=2.5cm]{geometry}
\usepackage{amsmath, amssymb, amsthm}
\usepackage{graphicx}
\usepackage{booktabs}
\usepackage{array}
\usepackage{multirow}
\usepackage{xcolor}
\usepackage{hyperref}
\usepackage{listings}
\usepackage{caption}
\usepackage{subcaption}
\usepackage{setspace}
\usepackage{fancyhdr}
\usepackage{titlesec}
\usepackage{abstract}
\usepackage{natbib}
\usepackage{url}

\usepackage{parskip}
\usepackage{tcolorbox}
\usepackage{algorithm}
\usepackage{algpseudocode}

% --- Colors ---
\definecolor{seoulblue}{RGB}{0, 71, 171}
\definecolor{codebg}{RGB}{245, 245, 245}
\definecolor{notecolor}{RGB}{255, 244, 206}
\definecolor{resultgreen}{RGB}{29, 158, 117}

% --- Hyperref setup ---
\hypersetup{
    colorlinks=true,
    linkcolor=seoulblue,
    citecolor=seoulblue,
    urlcolor=seoulblue,
    pdftitle={Neighbor-Interpolated Point Cloud Densification for 3D Gaussian Splatting},
    pdfauthor={Abduganiev Abdulvosit}
}

% --- Header/Footer ---
\pagestyle{fancy}
\fancyhf{}
\fancyhead[L]{\small\textit{Abduganiev Abdulvosit}}
\fancyhead[R]{\small\textit{MSc Thesis Draft, SeoulTech 2026}}
\fancyfoot[C]{\thepage}
\renewcommand{\headrulewidth}{0.4pt}

% --- Section formatting ---
\titleformat{\section}{\large\bfseries\color{seoulblue}}{{\thesection}}{1em}{}
\titleformat{\subsection}{\normalsize\bfseries}{{\thesubsection}}{1em}{}

% --- Code listing style ---
\lstset{
    backgroundcolor=\color{codebg},
    basicstyle=\ttfamily\footnotesize,
    breaklines=true,
    frame=single,
    rulecolor=\color{gray!40},
    keywordstyle=\color{seoulblue}\bfseries,
    commentstyle=\color{gray},
    stringstyle=\color{resultgreen},
    numbers=left,
    numberstyle=\tiny\color{gray},
    numbersep=6pt,
    xleftmargin=12pt,
    language=Python
}

% --- Theorem environments ---
\newtheorem{definition}{Definition}
\newtheorem{remark}{Remark}

% --- Note box ---
\tcbset{
    notebox/.style={
        colback=notecolor,
        colframe=orange!60,
        fonttitle=\bfseries,
        title=Note,
        left=4pt, right=4pt, top=4pt, bottom=4pt
    }
}

% ============================================================
\begin{document}
\onehalfspacing

% ---- TITLE PAGE ----
\begin{titlepage}
\centering
\vspace*{2cm}

{\Large\bfseries\color{seoulblue}
Neighbor-Interpolated Point Cloud Densification\\[0.5em]
for 3D Gaussian Splatting Initialization}

\vspace{1.5cm}
{\large Thesis Research Draft}

\vspace{2cm}
{\Large\bfseries ABDUGANIEV ABDULVOSIT}\\[0.5em]
{\normalsize Student ID: 24512169}

\vspace{1cm}
{\normalsize
Seoul National University of Science and Technology (SeoulTech)\\
Department of Integrated IT Engineering\\Smart Media and Network Lab
Master of Science\\
2026}

\vspace{1.5cm}
{\normalsize Supervisor: Prof. Dong Ho Kim}

\vspace{2cm}
{\normalsize
\textbf{GitHub:}
\url{https://github.com/Abdulbaasit98/3dgs-pointcloud-densification}}

\vfill
{\small Draft -- \today}
\end{titlepage}

% ---- TABLE OF CONTENTS ----
\tableofcontents
\newpage

% ============================================================
\section*{Abstract}
\addcontentsline{toc}{section}{Abstract}

3D Gaussian Splatting (3DGS) has emerged as a leading technique for real-time novel
view synthesis, achieving state-of-the-art quality through a differentiable Gaussian
rasterizer jointly optimized with a set of 3D Gaussians initialized from a sparse
Structure-from-Motion (SfM) point cloud. A known limitation of this pipeline is
sensitivity to the quality and coverage of the initial point cloud: regions
under-sampled by SfM---such as textureless or reflective surfaces---begin training
with insufficient Gaussian density and depend entirely on Adaptive Density Control
(ADC) to compensate, which can slow early convergence.

This thesis proposes and evaluates a lightweight preprocessing step---\textit{neighbor-interpolated point cloud densification}---that inserts new points at the
midpoints of geometrically sparse neighbor pairs in the SfM point cloud before
Gaussian initialization. The method requires zero modifications to the 3DGS training
code, runs on CPU in seconds, and increases the initial point count by a controlled
margin.

Experiments on the Tanks and Temples \textit{truck} scene demonstrate consistent
improvement in PSNR, SSIM, and LPIPS at every evaluated training checkpoint
($\Delta\text{PSNR} = +0.39\,\text{dB}$ at iteration 1{,}000;
$+0.17\,\text{dB}$ at 7{,}000), with the improvement most pronounced in early
training and diminishing as ADC compensates. Per-view statistical analysis
(Wilcoxon signed-rank, $p = 0.49$) indicates the effect, while consistent in
aggregate metrics, does not reach significance at the per-view level with $n = 32$
views and a single scene. Final Gaussian counts were comparable between conditions
(baseline: 1{,}681{,}300; densified: 1{,}733{,}031), indicating no compactness
benefit. All code, data, and results are publicly available.

\textbf{Keywords:} 3D Gaussian Splatting, Novel View Synthesis, Point Cloud
Initialization, Structure-from-Motion, Adaptive Density Control.

\newpage

% ============================================================
\section{Introduction}

Novel view synthesis---the problem of rendering a scene from arbitrary camera
positions given only a set of training images---has undergone a dramatic
transformation in recent years. Neural Radiance Fields (NeRF), introduced by
\citet{mildenhall2020nerf}, demonstrated that a continuous volumetric scene
representation could be learned end-to-end from posed images using differentiable
volume rendering, achieving unprecedented photorealism. Despite this success,
NeRF's reliance on ray marching through a neural network makes real-time rendering
computationally prohibitive.

3D Gaussian Splatting (3DGS), proposed by \citet{kerbl20233dgs}, addresses this
limitation by representing scenes as a collection of anisotropic 3D Gaussians
rendered through a differentiable tile-based rasterizer. By coupling an explicit
point-cloud representation with a fast GPU rasterizer, 3DGS achieves real-time
rendering speeds ($>$100 FPS) while matching or exceeding NeRF-level quality on
standard benchmarks.

A core aspect of the 3DGS pipeline that has received comparatively little attention
is the \textit{initialization} of the Gaussian set. The method initializes one
Gaussian at the location of each point in a sparse 3D point cloud produced by
Structure-from-Motion (SfM), typically COLMAP \citep{schonberger2016sfm}. Regions
that are geometrically challenging for SfM---including textureless surfaces,
specular materials, and weakly-observed scene boundaries---produce few or no
reconstructed points, resulting in gaps in the initial Gaussian set.

This thesis investigates a simple, targeted approach: gap-filling point cloud
densification via $k$-nearest-neighbor midpoint interpolation. Before the SfM
point cloud is passed to the Gaussian initializer, a lightweight preprocessing pass
identifies sparse neighbor pairs and inserts new points at their midpoints with
interpolated color. This step is purely geometric, runs entirely on CPU in seconds,
requires no neural components, and leaves the 3DGS training pipeline completely
unmodified.

\subsection{Contributions}

\begin{itemize}
    \item A gap-filling point cloud densification algorithm based on $k$-NN
    midpoint interpolation, implemented as a CPU-only preprocessing script
    compatible with the official 3DGS pipeline.

    \item A controlled experimental evaluation on the Tanks and Temples
    \textit{truck} scene comparing standard vs.\ densified initialization
    across multiple training checkpoints, using PSNR, SSIM, and LPIPS.

    \item A per-view robustness analysis including Wilcoxon signed-rank and
    paired $t$-test statistical testing, and a compactness analysis of
    final Gaussian counts under both conditions.

    \item A fully reproducible implementation with public code, notebook,
    and results at \url{https://github.com/Abdulbaasit98/3dgs-pointcloud-densification}.
\end{itemize}

% ============================================================
\section{Related Work}

\subsection{Neural Radiance Fields}

\citet{mildenhall2020nerf} introduced NeRF, which encodes a scene as a continuous
volumetric function $F_\theta: (\mathbf{x}, \mathbf{d}) \to (\mathbf{c}, \sigma)$
mapping a 3D position $\mathbf{x} \in \mathbb{R}^3$ and viewing direction
$\mathbf{d} \in \mathbb{S}^2$ to color and density, optimized via differentiable
volume rendering. Subsequent work improved training speed through hash-grid
encodings \citep{muller2022instant} and extended NeRF to unbounded scenes
\citep{barron2022mipnerf360} and dynamic scenes \citep{pumarola2021dnerf}.

\subsection{3D Gaussian Splatting}

\citet{kerbl20233dgs} proposed 3DGS, representing scenes as a set of anisotropic
3D Gaussians. Each Gaussian $\mathcal{G}_i$ is parameterized by:
\begin{equation}
    \mathcal{G}_i = \{\boldsymbol{\mu}_i,\, \boldsymbol{\Sigma}_i,\, \alpha_i,\, \mathbf{c}_i\}
\end{equation}
where $\boldsymbol{\mu}_i \in \mathbb{R}^3$ is the center position,
$\boldsymbol{\Sigma}_i \in \mathbb{R}^{3\times3}$ is the covariance matrix
(decomposed as $\boldsymbol{\Sigma} = \mathbf{R}\mathbf{S}\mathbf{S}^\top\mathbf{R}^\top$
with rotation $\mathbf{R}$ and scaling $\mathbf{S}$), $\alpha_i \in [0,1]$ is
opacity, and $\mathbf{c}_i$ are spherical harmonic color coefficients.

The influence of a Gaussian at a 3D point $\mathbf{x}$ is:
\begin{equation}
    G_i(\mathbf{x}) = \exp\!\left(-\tfrac{1}{2}(\mathbf{x}-\boldsymbol{\mu}_i)^\top \boldsymbol{\Sigma}_i^{-1} (\mathbf{x}-\boldsymbol{\mu}_i)\right)
\end{equation}

Rendering proceeds via $\alpha$-compositing of Gaussians sorted by depth:
\begin{equation}
    \mathbf{C} = \sum_{i=1}^{N} \mathbf{c}_i \alpha_i \prod_{j=1}^{i-1}(1 - \alpha_j)
\end{equation}

Optimization uses a combined photometric and SSIM loss:
\begin{equation}
    \mathcal{L} = (1 - \lambda)\mathcal{L}_1 + \lambda\,\mathcal{L}_\text{SSIM}
    \label{eq:loss}
\end{equation}
where $\lambda = 0.2$ in the original paper.

\subsection{Initialization Sensitivity in 3DGS}

\citet{foroutan2024sfm} systematically evaluated alternatives to SfM initialization
for Gaussian Splatting, finding that initialization quality significantly affects
convergence and final reconstruction quality. \citet{jung2024rain} proposed RAIN-GS,
which relaxes the accurate initialization constraint, allowing training to recover
from low-quality point clouds through aggressive optimization. The present work
takes the complementary direction: improving the initialization itself via fast
geometric preprocessing, leaving the optimizer unchanged.

\subsection{Point Cloud Densification}

Classical point cloud upsampling methods include Moving Least Squares (MLS)
surface fitting and deep learning approaches such as PU-Net. These target
production of geometrically accurate dense surfaces. Our approach is simpler
and more targeted: reducing large gaps in the SfM cloud to improve Gaussian
coverage, using only pairwise distance geometry and linear color interpolation---
no surface model is required.

\subsection{Compact Gaussian Representations}

LightGaussian \citep{fan2024lightgaussian} and Compact3D \citep{navaneet2023compact3d}
propose pruning and quantization strategies to reduce the final Gaussian count.
Scaffold-GS \citep{lu2024scaffold} introduces anchor-based Gaussians to reduce
redundancy. These methods operate on the trained model, complementary to the
present work. The compactness analysis in Section~\ref{sec:compactness} directly
examines whether initialization density affects the final Gaussian count.

% ============================================================
\section{Method}

\subsection{Problem Formulation}

Let $\mathcal{P} = \{(\mathbf{p}_i, \mathbf{c}_i)\}_{i=1}^{N}$ denote the sparse
SfM point cloud, where $\mathbf{p}_i \in \mathbb{R}^3$ is the 3D position and
$\mathbf{c}_i \in [0,255]^3$ is the RGB color of the $i$-th point. The 3DGS
pipeline initializes one Gaussian at each $\mathbf{p}_i$. We seek to construct a
denser point cloud
$\mathcal{P}' = \mathcal{P} \cup \mathcal{P}_\text{new}$
such that the geometric coverage of the scene is improved, without access to
additional images or depth sensors.

\subsection{Gap Detection via $k$-Nearest Neighbors}

For each point $\mathbf{p}_i \in \mathcal{P}$, we compute its $k$ nearest
neighbors using a $k$-d tree. Let $d_{ij} = \|\mathbf{p}_i - \mathbf{p}_j\|_2$
denote the Euclidean distance from point $i$ to its $j$-th nearest neighbor.

We define the \textit{scene-scale reference distance} as the median nearest-neighbor
distance:
\begin{equation}
    d_\text{med} = \mathrm{median}_{i}\, d_{i,1}
    \label{eq:dmed}
\end{equation}

A neighbor pair $(\mathbf{p}_i, \mathbf{p}_j)$ is classified as a \textit{gap} if:
\begin{equation}
    d_{ij} > \tau \cdot d_\text{med}
    \label{eq:gap}
\end{equation}
where $\tau > 1$ is the gap-threshold factor (hyperparameter). A higher $\tau$
results in fewer, more conservative insertions; a lower $\tau$ fills more aggressively.

\subsection{Midpoint Interpolation}

For each detected gap pair $(\mathbf{p}_i, \mathbf{p}_j)$, a new point is inserted
at their geometric midpoint with linearly interpolated color:
\begin{align}
    \mathbf{p}_\text{new} &= \frac{\mathbf{p}_i + \mathbf{p}_j}{2}
    \label{eq:midpoint} \\
    \mathbf{c}_\text{new} &= \left\lfloor\frac{\mathbf{c}_i + \mathbf{c}_j}{2}\right\rfloor
    \label{eq:colorinterp}
\end{align}

To prevent excessive growth in sparse outlier regions, the number of new points
inserted per original point is capped at $m_\text{max}$ (hyperparameter). Each
pair $(i, j)$ is processed at most once (using a seen-pairs set).

\subsection{Algorithm}

\begin{algorithm}[H]
\caption{Neighbor-Interpolated Point Cloud Densification}
\label{alg:densify}
\begin{algorithmic}[1]
\Require Point cloud $\mathcal{P} = \{(\mathbf{p}_i, \mathbf{c}_i)\}$,
         neighbors $k$, threshold factor $\tau$, max new per point $m_\text{max}$
\Ensure Densified point cloud $\mathcal{P}'$
\State Build $k$-d tree $T$ over $\{\mathbf{p}_i\}$
\State Compute $d_\text{med}$ via Eq.~\eqref{eq:dmed}
\State $\mathcal{P}_\text{new} \leftarrow \emptyset$; $\text{seen} \leftarrow \emptyset$
\For{each point $i$}
    \State $\text{inserted} \leftarrow 0$
    \For{each neighbor $j$ of $i$ in $T$ (in distance order)}
        \If{$\text{inserted} \geq m_\text{max}$} \textbf{break} \EndIf
        \If{$d_{ij} \leq \tau \cdot d_\text{med}$} \textbf{continue} \EndIf
        \If{$(i,j) \in \text{seen}$} \textbf{continue} \EndIf
        \State $\text{seen} \leftarrow \text{seen} \cup \{(\min(i,j),\max(i,j))\}$
        \State Compute $\mathbf{p}_\text{new}$, $\mathbf{c}_\text{new}$ via Eqs.~\eqref{eq:midpoint}--\eqref{eq:colorinterp}
        \State $\mathcal{P}_\text{new} \leftarrow \mathcal{P}_\text{new} \cup \{(\mathbf{p}_\text{new}, \mathbf{c}_\text{new})\}$
        \State $\text{inserted} \leftarrow \text{inserted} + 1$
    \EndFor
\EndFor
\State \Return $\mathcal{P}' = \mathcal{P} \cup \mathcal{P}_\text{new}$
\end{algorithmic}
\end{algorithm}

\subsection{Hyperparameters}

Table~\ref{tab:hyperparams} summarizes the hyperparameters used in this study.

\begin{table}[h]
\centering
\caption{Hyperparameter settings used in the experiment.}
\label{tab:hyperparams}
\begin{tabular}{llp{7cm}}
\toprule
\textbf{Parameter} & \textbf{Value} & \textbf{Effect} \\
\midrule
$k$ & 6 & Number of nearest neighbors considered per point \\
$\tau$ & 3.0 & Gap threshold multiplier; higher = fewer insertions \\
$m_\text{max}$ & 1 & Max new points per original point; limits growth \\
\bottomrule
\end{tabular}
\end{table}

\subsection{Integration with 3DGS}

The method requires \textbf{zero changes} to the 3DGS training code. The
\texttt{dataset\_readers.py} loader in the official repository reads whichever
\texttt{points3D.ply} file it finds in the scene's \texttt{sparse/0/} folder.
Integration is therefore a file replacement:

\begin{enumerate}
    \item Copy the original scene folder (preserving cameras and images unchanged).
    \item Run Algorithm~\ref{alg:densify} on the original \texttt{points3D.bin},
          writing the output as \texttt{points3D.ply}.
    \item Train using the standard \texttt{train.py} command, pointing to the
          modified scene folder.
\end{enumerate}

% ============================================================
\section{Experiments}

\subsection{Scene and Dataset}

We evaluate on the \textit{truck} scene from the Tanks and Temples benchmark
\citep{knapitsch2017tanks}: 251 training images, 32 held-out test views,
approximately 136{,}029 points in the original SfM reconstruction.

\subsection{Training Configuration}

\begin{table}[h]
\centering
\caption{Training configuration for both conditions.}
\label{tab:config}
\begin{tabular}{ll}
\toprule
\textbf{Setting} & \textbf{Value} \\
\midrule
GPU & NVIDIA Tesla T4 (16\,GB, Google Colab free tier) \\
Iterations & 7{,}000 (paper uses 30{,}000; compute-constrained) \\
Test iterations & 1{,}000 / 3{,}000 / 7{,}000 \\
Loss & $(1-\lambda)\mathcal{L}_1 + \lambda\mathcal{L}_\text{SSIM}$, $\lambda=0.2$ \\
Random seed & Fixed (identical for both conditions) \\
\bottomrule
\end{tabular}
\end{table}

\subsection{Evaluation Metrics}

We report three standard metrics on held-out test views:

\begin{itemize}
    \item \textbf{PSNR} (Peak Signal-to-Noise Ratio, dB): higher is better.
    \begin{equation}
        \text{PSNR} = 10\log_{10}\!\left(\frac{\text{MAX}^2}{\text{MSE}}\right)
    \end{equation}
    \item \textbf{SSIM} (Structural Similarity Index): higher is better, $\in [0,1]$.
    \item \textbf{LPIPS} (Learned Perceptual Image Patch Similarity,
    VGG backbone): lower is better.
\end{itemize}

% ============================================================
\section{Results}

\subsection{Point Cloud Statistics}

After applying Algorithm~\ref{alg:densify} with $k=6$, $\tau=3.0$, $m_\text{max}=1$:

\begin{equation}
    d_\text{med} = 0.01129, \qquad \tau \cdot d_\text{med} = 0.03387
\end{equation}

\begin{table}[h]
\centering
\caption{Point cloud statistics before and after densification.}
\label{tab:pointcloud}
\begin{tabular}{lrr}
\toprule
\textbf{Condition} & \textbf{Initial points} & \textbf{Increase} \\
\midrule
Baseline (original SfM) & 136{,}029 & --- \\
Densified               & 197{,}869 & $+61{,}840$ ($+45.5\%$) \\
\bottomrule
\end{tabular}
\end{table}

\subsection{Convergence Speed}
\label{sec:convergence}

Table~\ref{tab:results} reports PSNR, SSIM, and LPIPS on 32 held-out test views at
matched training checkpoints. $\Delta$ values are densified minus baseline; positive
$\Delta$PSNR/$\Delta$SSIM and negative $\Delta$LPIPS indicate improvement.

\begin{table}[h]
\centering
\caption{Quantitative results: baseline vs.\ densified initialization.
Arrows indicate the direction of improvement.
$\Delta$ = densified $-$ baseline.}
\label{tab:results}
\begin{tabular}{c
    r
    r
    r
    r
    r
    r
    r
    r
    r}
\toprule
& \multicolumn{3}{c}{\textbf{PSNR} (dB) $\uparrow$}
& \multicolumn{3}{c}{\textbf{SSIM} $\uparrow$}
& \multicolumn{3}{c}{\textbf{LPIPS} $\downarrow$} \\
\cmidrule(lr){2-4}\cmidrule(lr){5-7}\cmidrule(lr){8-10}
\textbf{Iter} & {Base} & {Dense} & {$\Delta$}
              & {Base} & {Dense} & {$\Delta$}
              & {Base} & {Dense} & {$\Delta$} \\
\midrule
1{,}000 & 20.978 & 21.370 & \bfseries +0.391
        & 0.7315 & 0.7538 & \bfseries +0.0223
        & 0.3542 & 0.3264 & \bfseries -0.0278 \\
3{,}000 & 22.728 & 22.966 & \bfseries +0.238
        & 0.8074 & 0.8183 & \bfseries +0.0108
        & 0.2524 & 0.2368 & \bfseries -0.0156 \\
7{,}000 & 23.988 & 24.154 & \bfseries +0.165
        & 0.8528 & 0.8595 & \bfseries +0.0067
        & 0.1950 & 0.1846 & \bfseries -0.0103 \\
\midrule
\multicolumn{10}{l}{\textit{Reference (baseline only, full 30k run):}} \\
30{,}000 & 25.481 & {---} & {---}
         & 0.8850 & {---} & {---}
         & 0.1418 & {---} & {---} \\
\bottomrule
\end{tabular}
\end{table}

Densified initialization outperforms the baseline on all three metrics at every
checkpoint. The advantage is largest at iteration 1{,}000 ($\Delta\text{PSNR} =
+0.391\,\text{dB}$, $\Delta\text{SSIM} = +0.022$, $\Delta\text{LPIPS} = -0.028$)
and narrows monotonically through iteration 7{,}000 ($\Delta\text{PSNR} =
+0.165\,\text{dB}$). This convergence-speed pattern is consistent with the
hypothesis that denser initialization provides an early-training head start that
ADC gradually compensates for.

\subsection{Per-View Robustness Analysis}
\label{sec:perviewstats}

To assess whether the aggregate improvement held consistently across test views,
per-pixel $\ell_1$ error was computed individually for each of the $n=32$ test
views at iteration 7{,}000. Let $e_i^\text{base}$ and $e_i^\text{dense}$ denote
the mean $\ell_1$ error for view $i$; define $\delta_i = e_i^\text{dense} -
e_i^\text{base}$ (negative = densified better).

\begin{table}[h]
\centering
\caption{Per-view $\ell_1$ error analysis at iteration 7{,}000.}
\label{tab:perviewstats}
\begin{tabular}{lc}
\toprule
\textbf{Statistic} & \textbf{Value} \\
\midrule
Views where densified is better ($\delta_i < 0$) & $17/32$ ($53.1\%$) \\
Views where baseline is better ($\delta_i > 0$) & $15/32$ ($46.9\%$) \\
Mean delta $\bar{\delta}$ & $-0.037$ \\
Standard deviation $\sigma_\delta$ & $0.326$ \\
Best case (view 00014) & $\delta = -0.82$ \\
Worst case (view 00003) & $\delta = +0.68$ \\
\midrule
Paired $t$-test: $t$-statistic & $-0.634$ \\
Paired $t$-test: $p$-value & $0.531$ \\
Wilcoxon signed-rank: $p$-value & $0.488$ \\
\bottomrule
\end{tabular}
\end{table}

Both the paired $t$-test ($p = 0.531$) and the Wilcoxon signed-rank test
($p = 0.488$) fail to reject the null hypothesis ($\alpha = 0.05$). The aggregate
metric improvement, while consistent and reproducible, is not statistically
significant at the per-view level given $n = 32$ views and a single scene/seed.

\subsection{Compactness Analysis}
\label{sec:compactness}

Table~\ref{tab:compactness} reports the final Gaussian count for both conditions
at iteration 7{,}000.

\begin{table}[h]
\centering
\caption{Final Gaussian counts at iteration 7{,}000.}
\label{tab:compactness}
\begin{tabular}{lrrrr}
\toprule
\textbf{Condition} & \textbf{Init points} & \textbf{Final Gaussians}
                   & \textbf{ADC added} & \textbf{Increase} \\
\midrule
Baseline  & 136{,}029 & 1{,}681{,}300 & $+1{,}545{,}271$ & $12.4\times$ \\
Densified & 197{,}869 & 1{,}733{,}031 & $+1{,}535{,}162$ & $\phantom{1}8.8\times$ \\
\midrule
$\Delta$ & $+45.5\%$ & $+3.1\%$ & & \\
\bottomrule
\end{tabular}
\end{table}

Despite a 45.5\% larger initialization, the densified model contains only 3.1\%
more Gaussians at convergence. ADC added approximately $1.54\times10^6$ Gaussians
in both conditions, suggesting its growth criteria respond to per-view residual
error rather than global point density. The proposed method therefore confers
\textbf{no compactness benefit}.

\subsection{Summary}
\label{sec:summary}

\begin{tcolorbox}[notebox]
Denser initialization produces a small, consistent improvement in aggregate
reconstruction quality during early training ($\Delta\text{PSNR} = +0.39\,\text{dB}$
at iter.\ 1{,}000). This effect is not statistically significant at the per-view
level ($p \approx 0.49$, $n = 32$) and does not reduce the final Gaussian count.
The contribution is best characterized as a modest, directional early-convergence
benefit rather than a definitive quality improvement.
\end{tcolorbox}

% ============================================================
\section{Limitations}

\begin{itemize}
    \item \textbf{Single scene:} Evaluation is limited to the Tanks and Temples
    \textit{truck} scene. Generalization to other scene types is unknown.

    \item \textbf{Single seed:} All runs use a fixed random seed. Results may
    vary across seeds; the per-view analysis ($p \approx 0.49$) is inconclusive
    with $n = 32$.

    \item \textbf{Reduced iterations:} Training was capped at 7{,}000 iterations
    due to Google Colab free-tier session limits. The reference paper uses 30{,}000
    iterations. The baseline's 30{,}000-iteration result is provided in
    Table~\ref{tab:results} for reference only.

    \item \textbf{No hyperparameter sweep:} Only one setting of $(\tau, k, m_\text{max})$
    was evaluated due to compute constraints.

    \item \textbf{No surface geometry:} The midpoint interpolation in
    Eq.~\eqref{eq:midpoint} does not account for surface orientation or local
    geometry, which may lead to off-surface insertions in curved regions.
\end{itemize}

% ============================================================
\section{Conclusion}

This thesis proposed and evaluated neighbor-interpolated point cloud densification
as a preprocessing strategy for improving 3D Gaussian Splatting initialization.
The core contribution is a simple geometric algorithm (Algorithm~\ref{alg:densify})
that identifies and fills gaps in sparse SfM point clouds using $k$-NN midpoint
interpolation, requiring no changes to the 3DGS training pipeline.

Experiments on the Tanks and Temples \textit{truck} scene demonstrated a consistent
aggregate improvement across PSNR, SSIM, and LPIPS at every evaluated checkpoint
(Table~\ref{tab:results}), with the largest advantage early in training
($\Delta\text{PSNR} = +0.391\,\text{dB}$ at iteration 1{,}000) diminishing toward
convergence. Two important null findings were also reported: (1) the improvement is
not statistically significant at the per-view level (Wilcoxon $p = 0.488$,
Table~\ref{tab:perviewstats}); and (2) the method provides no compactness benefit,
as Adaptive Density Control operates independently of initialization density
(Table~\ref{tab:compactness}).

These findings position the method as a lightweight, zero-cost-at-training-time
initialization improvement with a modest, directional benefit --- a complete,
honest, reproducible MSc-level experiment with clear directions for future work.

\subsection*{Future Work}

\begin{itemize}
    \item Multi-scene validation across the full Tanks and Temples and Deep
    Blending benchmarks to assess generalization.

    \item Hyperparameter sweep over $\tau \in \{2, 3, 4, 5\}$ and
    $k \in \{4, 6, 8\}$ to characterize the quality/compactness trade-off.

    \item Full 30{,}000-iteration runs to determine whether the early-training
    advantage persists at convergence.

    \item Surface-aware insertions using local normal estimation to improve
    geometric accuracy of synthetic points.

    \item Extension toward an IEEE Access submission with multi-scene,
    multi-seed statistical validation.
\end{itemize}

% ============================================================
\bibliographystyle{plainnat}
\begin{thebibliography}{99}

\bibitem[Barron et al.(2022)]{barron2022mipnerf360}
Barron, J.~T., Mildenhall, B., Verbin, D., Srinivasan, P.~P., and Hedman, P.
\newblock Mip-{NeRF} 360: Unbounded anti-aliased neural radiance fields.
\newblock In \emph{CVPR}, 2022.

\bibitem[Fan et al.(2024)]{fan2024lightgaussian}
Fan, Z. et al.
\newblock {LightGaussian}: Unbounded 3{D} {Gaussian} compression with 15x
reduction and 200+ {FPS}.
\newblock In \emph{NeurIPS}, 2024.

\bibitem[Foroutan et al.(2024)]{foroutan2024sfm}
Foroutan, Y., Rebain, D., Yi, K.~M., and Tagliasacchi, A.
\newblock Evaluating alternatives to {SfM} point cloud initialization for
{Gaussian} splatting.
\newblock \emph{arXiv:2404.12547}, 2024.

\bibitem[Jung et al.(2024)]{jung2024rain}
Jung, J., Han, J., An, H., Kang, J., Park, S., and Kim, S.
\newblock Relaxing accurate initialization constraint for {3D} {Gaussian}
splatting.
\newblock \emph{arXiv:2403.09413}, 2024.

\bibitem[Kerbl et al.(2023)]{kerbl20233dgs}
Kerbl, B., Kopanas, G., Leimk\"{u}hler, T., and Drettakis, G.
\newblock {3D Gaussian} splatting for real-time radiance field rendering.
\newblock \emph{ACM Transactions on Graphics}, 42(4), 2023.

\bibitem[Knapitsch et al.(2017)]{knapitsch2017tanks}
Knapitsch, A., Park, J., Zhou, Q.-Y., and Koltun, V.
\newblock Tanks and temples: Benchmarking large-scale scene reconstruction.
\newblock \emph{ACM Transactions on Graphics}, 36(4), 2017.

\bibitem[Lu et al.(2024)]{lu2024scaffold}
Lu, T., Yu, M., Xu, L., Xiangli, Y., Wang, L., Lin, D., and Dai, B.
\newblock {Scaffold-GS}: Structured 3{D} {Gaussians} for view-adaptive
rendering.
\newblock In \emph{CVPR}, 2024.

\bibitem[Mildenhall et al.(2020)]{mildenhall2020nerf}
Mildenhall, B., Srinivasan, P.~P., Tancik, M., Barron, J.~T., Ramamoorthi, R.,
and Ng, R.
\newblock {NeRF}: Representing scenes as neural radiance fields for view
synthesis.
\newblock In \emph{ECCV}, 2020.

\bibitem[M\"{u}ller et al.(2022)]{muller2022instant}
M\"{u}ller, T., Evans, A., Schied, C., and Keller, A.
\newblock Instant neural graphics primitives with a multiresolution hash
encoding.
\newblock \emph{ACM Transactions on Graphics}, 41(4), 2022.

\bibitem[Navaneet et al.(2023)]{navaneet2023compact3d}
Navaneet, K.~L., Meibodi, K.~P., Koohpayegani, S.~A., and Pirsiavash, H.
\newblock {Compact3D}: Compressing {Gaussian} splat radiance field models with
vector quantization.
\newblock \emph{arXiv:2311.18159}, 2023.

\bibitem[Pumarola et al.(2021)]{pumarola2021dnerf}
Pumarola, A., Corona, E., Pons-Moll, G., and Moreno-Noguer, F.
\newblock {D-NeRF}: Neural radiance fields for dynamic scenes.
\newblock In \emph{CVPR}, 2021.

\bibitem[Sch\"{o}nberger and Frahm(2016)]{schonberger2016sfm}
Sch\"{o}nberger, J.~L. and Frahm, J.-M.
\newblock Structure-from-motion revisited.
\newblock In \emph{CVPR}, 2016.

\end{thebibliography}

\end{document}
